\documentclass[11pt,a4paper]{article}

\usepackage[utf8]{inputenc}
\usepackage[T1]{fontenc}
\usepackage[french]{babel}
\usepackage[a4paper,margin=2.2cm]{geometry}
\usepackage{amsmath,amssymb,amsthm,mathtools}
\usepackage{lmodern}
\usepackage{enumitem}
\usepackage{hyperref}
\usepackage{xcolor}
\usepackage{fancyhdr}
\usepackage{stmaryrd}
\usepackage{mathrsfs}
\usepackage{array}

\hypersetup{
    colorlinks=true,
    linkcolor=black,
    urlcolor=blue,
    citecolor=black
}

\setlength{\parindent}{0pt}
\setlength{\parskip}{0.6em}

\pagestyle{fancy}
\fancyhf{}
\fancyhead[L]{Algèbre linéaire --- Chapitre 9}
\fancyhead[R]{Rang}
\fancyfoot[C]{\thepage}

% Environnements
\newtheorem{definition}{Définition}[section]
\newtheorem{theoreme}[definition]{Théorème}
\newtheorem{proposition}[definition]{Proposition}
\newtheorem{corollaire}[definition]{Corollaire}
\newtheorem{lemme}[definition]{Lemme}
\theoremstyle{remark}
\newtheorem{remarque}[definition]{Remarque}
\newtheorem{exemple}[definition]{Exemple}
\newtheorem*{methode}{Méthode}
\newtheorem*{idee}{Idée}

% Commandes
\newcommand{\R}{\mathbb{R}}
\newcommand{\C}{\mathbb{C}}
\newcommand{\K}{\mathbb{K}}
\newcommand{\N}{\mathbb{N}}
\newcommand{\Z}{\mathbb{Z}}
\newcommand{\Vect}{\mathrm{Vect}}
\newcommand{\Ker}{\mathrm{Ker}}
\newcommand{\Img}{\mathrm{Im}}
\newcommand{\rg}{\mathrm{rg}}
\newcommand{\id}{\mathrm{Id}}
\newcommand{\Lin}{\mathcal{L}}
\newcommand{\Mnp}{M_{n,p}(\K)}
\newcommand{\Mpn}{M_{p,n}(\K)}
\newcommand{\Mnn}{M_n(\K)}
\newcommand{\GLn}{GL_n(\K)}

\begin{document}

\begin{center}
    {\LARGE \textbf{Chapitre 9 --- Rang}}\\[0.4em]
    {\large Version complète et détaillée}
\end{center}

\hrule
\vspace{1em}

\tableofcontents

\newpage

\section{Pourquoi la notion de rang est-elle fondamentale ?}

En algèbre linéaire, on rencontre sans cesse des objets qui contiennent de la redondance :
\begin{itemize}
    \item une famille de vecteurs peut contenir des vecteurs inutiles ;
    \item une matrice peut avoir plusieurs colonnes dépendantes ;
    \item une application linéaire peut envoyer tout un espace dans un sous-espace beaucoup plus petit.
\end{itemize}

La notion de \textbf{rang} sert à mesurer la quantité de structure linéaire réellement produite.

\begin{idee}
Le rang mesure le nombre de directions indépendantes effectivement présentes.
\end{idee}

Le rang apparaît naturellement :
\begin{itemize}
    \item lorsqu’on étudie une famille de vecteurs ;
    \item lorsqu’on étudie l’image d’une application linéaire ;
    \item lorsqu’on travaille avec les colonnes ou les lignes d’une matrice ;
    \item lorsqu’on résout des systèmes linéaires.
\end{itemize}

\begin{remarque}
C’est une notion unificatrice : derrière le rang d’une famille, le rang d’une matrice et le rang d’une application linéaire, il y a exactement la même idée.
\end{remarque}

\section{Rang d’une famille de vecteurs}

\subsection{Définition}

\begin{definition}
Soit \((u_1,\dots,u_p)\) une famille finie de vecteurs d’un espace vectoriel \(E\).  
On appelle \textbf{rang} de cette famille la dimension du sous-espace qu’elle engendre :
\[
\rg(u_1,\dots,u_p)=\dim\big(\Vect(u_1,\dots,u_p)\big).
\]
\end{definition}

\begin{remarque}
Le rang ne dépend donc pas du nombre de vecteurs de départ, mais seulement du sous-espace qu’ils engendrent effectivement.
\end{remarque}

\begin{exemple}
Dans \(\R^3\), considérons
\[
u_1=(1,0,0),\qquad u_2=(0,1,0),\qquad u_3=(1,1,0).
\]
Comme
\[
u_3=u_1+u_2,
\]
on a
\[
\Vect(u_1,u_2,u_3)=\Vect(u_1,u_2),
\]
donc
\[
\rg(u_1,u_2,u_3)=2.
\]
\end{exemple}

\subsection{Premières propriétés}

\begin{proposition}
Pour toute famille finie \((u_1,\dots,u_p)\),
\[
0\le \rg(u_1,\dots,u_p)\le p.
\]
\end{proposition}

\begin{proof}
Le rang est une dimension, donc il est positif ou nul.

D’autre part, l’espace \(\Vect(u_1,\dots,u_p)\) est engendré par \(p\) vecteurs. Sa dimension est donc inférieure ou égale à \(p\).
\end{proof}

\begin{proposition}
La famille \((u_1,\dots,u_p)\) est libre si et seulement si
\[
\rg(u_1,\dots,u_p)=p.
\]
\end{proposition}

\begin{proof}
La famille est libre si et seulement si elle forme une base de son espace engendré. Cela signifie exactement que la dimension de cet espace est égale au nombre de vecteurs, soit \(p\).
\end{proof}

\begin{proposition}
Si la famille \((u_1,\dots,u_p)\) engendre un espace \(F\), alors
\[
\rg(u_1,\dots,u_p)=\dim(F).
\]
\end{proposition}

\begin{proof}
Si la famille engendre \(F\), alors
\[
\Vect(u_1,\dots,u_p)=F.
\]
On prend la dimension.
\end{proof}

\begin{corollaire}
Une famille est génératrice de \(E\) si et seulement si
\[
\rg(u_1,\dots,u_p)=\dim(E).
\]
\end{corollaire}

\subsection{Interprétation par extraction de base}

\begin{proposition}
Le rang d’une famille est le nombre de vecteurs d’une base extraite de cette famille.
\end{proposition}

\begin{proof}
Le sous-espace
\[
F=\Vect(u_1,\dots,u_p)
\]
est engendré par la famille. On peut en extraire une base. Le nombre de vecteurs de cette base vaut
\[
\dim(F)=\rg(u_1,\dots,u_p).
\]
\end{proof}

\begin{remarque}
On peut donc voir le rang comme :
\begin{center}
\textit{le nombre maximal de vecteurs linéairement indépendants que l’on peut extraire de la famille.}
\end{center}
\end{remarque}

\section{Rang d’une application linéaire}

\subsection{Définition}

\begin{definition}
Soit
\[
u:E\to F
\]
une application linéaire, avec \(E\) de dimension finie.  
On appelle \textbf{rang} de \(u\), noté \(\rg(u)\), la dimension de son image :
\[
\rg(u)=\dim(\Img(u)).
\]
\end{definition}

\begin{exemple}
Soit
\[
u:\R^3\to\R^2,\qquad u(x,y,z)=(x+y,\ y+z).
\]
Alors
\[
u(1,0,0)=(1,0),\qquad u(0,1,0)=(1,1),\qquad u(0,0,1)=(0,1).
\]
Les vecteurs \((1,0)\) et \((0,1)\) appartiennent à l’image, donc
\[
\Img(u)=\R^2.
\]
Ainsi
\[
\rg(u)=2.
\]
\end{exemple}

\subsection{Premières bornes}

\begin{proposition}
Si \(u:E\to F\) est linéaire entre espaces de dimension finie, alors
\[
0\le \rg(u)\le \dim(F)
\qquad \text{et} \qquad
\rg(u)\le \dim(E).
\]
\end{proposition}

\begin{proof}
Comme \(\Img(u)\) est un sous-espace de \(F\),
\[
\dim(\Img(u))\le \dim(F).
\]
Donc
\[
\rg(u)\le \dim(F).
\]
D’autre part, le théorème du rang donnera
\[
\dim(E)=\dim(\Ker(u))+\rg(u),
\]
d’où
\[
\rg(u)\le \dim(E).
\]
\end{proof}

\begin{corollaire}
On a toujours
\[
\rg(u)\le \min(\dim(E),\dim(F)).
\]
\end{corollaire}

\section{Théorème du rang}

\subsection{Énoncé}

\begin{theoreme}[Théorème du rang]
Soit
\[
u:E\to F
\]
une application linéaire, avec \(E\) de dimension finie. Alors
\[
\dim(E)=\dim(\Ker(u))+\rg(u).
\]
\end{theoreme}

\subsection{Démonstration}

\begin{proof}
Posons
\[
k=\dim(\Ker(u)).
\]
Choisissons une base du noyau :
\[
(e_1,\dots,e_k).
\]
Comme \(E\) est de dimension finie, on peut compléter cette famille en une base de \(E\) :
\[
(e_1,\dots,e_k,e_{k+1},\dots,e_n),
\]
où
\[
n=\dim(E).
\]

Nous allons montrer que la famille
\[
\big(u(e_{k+1}),\dots,u(e_n)\big)
\]
est une base de \(\Img(u)\).

\textbf{1. Elle engendre \(\Img(u)\).}

Soit \(y\in \Img(u)\). Il existe \(x\in E\) tel que
\[
y=u(x).
\]
Comme \((e_1,\dots,e_n)\) est une base de \(E\), on peut écrire
\[
x=\lambda_1e_1+\cdots+\lambda_ne_n.
\]
Alors
\[
y=u(x)=\lambda_1u(e_1)+\cdots+\lambda_nu(e_n).
\]
Mais pour \(i\le k\), comme \(e_i\in \Ker(u)\),
\[
u(e_i)=0.
\]
Donc
\[
y=\lambda_{k+1}u(e_{k+1})+\cdots+\lambda_nu(e_n).
\]
Ainsi, la famille engendre \(\Img(u)\).

\textbf{2. Elle est libre.}

Supposons
\[
\lambda_{k+1}u(e_{k+1})+\cdots+\lambda_nu(e_n)=0.
\]
Alors
\[
u(\lambda_{k+1}e_{k+1}+\cdots+\lambda_ne_n)=0,
\]
donc
\[
\lambda_{k+1}e_{k+1}+\cdots+\lambda_ne_n\in \Ker(u).
\]
Or
\[
\Ker(u)=\Vect(e_1,\dots,e_k).
\]
Donc il existe \(\mu_1,\dots,\mu_k\) tels que
\[
\lambda_{k+1}e_{k+1}+\cdots+\lambda_ne_n
=
\mu_1e_1+\cdots+\mu_ke_k.
\]
En réorganisant,
\[
-\mu_1e_1-\cdots-\mu_ke_k+\lambda_{k+1}e_{k+1}+\cdots+\lambda_ne_n=0.
\]
Comme \((e_1,\dots,e_n)\) est une base, elle est libre. Tous les coefficients sont donc nuls :
\[
\mu_1=\cdots=\mu_k=\lambda_{k+1}=\cdots=\lambda_n=0.
\]
La famille est libre.

On a donc montré que
\[
\big(u(e_{k+1}),\dots,u(e_n)\big)
\]
est une base de \(\Img(u)\). Elle contient
\[
n-k
\]
vecteurs. Ainsi
\[
\rg(u)=n-k.
\]
Comme \(n=\dim(E)\) et \(k=\dim(\Ker(u))\), on obtient
\[
\dim(E)=\dim(\Ker(u))+\rg(u).
\]
\end{proof}

\subsection{Conséquences fondamentales}

\begin{corollaire}
Soit \(u:E\to F\) linéaire, avec \(E\) de dimension finie. Alors
\[
u \text{ injective }
\Longleftrightarrow
\rg(u)=\dim(E).
\]
\end{corollaire}

\begin{proof}
L’injectivité équivaut à
\[
\Ker(u)=\{0\},
\]
donc à
\[
\dim(\Ker(u))=0.
\]
Le théorème du rang donne alors
\[
\dim(E)=\rg(u).
\]
\end{proof}

\begin{corollaire}
Soit \(u:E\to F\) linéaire, avec \(F\) de dimension finie. Alors
\[
u \text{ surjective }
\Longleftrightarrow
\rg(u)=\dim(F).
\]
\end{corollaire}

\begin{proof}
La surjectivité équivaut à
\[
\Img(u)=F,
\]
donc à
\[
\dim(\Img(u))=\dim(F).
\]
\end{proof}

\begin{corollaire}
Si
\[
\dim(E)=\dim(F)<+\infty,
\]
alors pour une application linéaire \(u:E\to F\), les propriétés suivantes sont équivalentes :
\begin{itemize}
    \item \(u\) est injective ;
    \item \(u\) est surjective ;
    \item \(u\) est bijective ;
    \item \(\rg(u)=\dim(E)=\dim(F)\).
\end{itemize}
\end{corollaire}

\begin{proof}
Si \(u\) est injective, alors
\[
\rg(u)=\dim(E)=\dim(F),
\]
donc \(u\) est surjective.

Si \(u\) est surjective, alors
\[
\rg(u)=\dim(F)=\dim(E),
\]
donc \(\Ker(u)=\{0\}\), et \(u\) est injective.
\end{proof}

\section{Rang d’une matrice}

\subsection{Définition}

\begin{definition}
Soit
\[
A\in M_{p,n}(\K).
\]
On lui associe l’application linéaire
\[
u_A:\K^n\to\K^p,\qquad X\mapsto AX.
\]
On appelle \textbf{rang de la matrice \(A\)} le rang de \(u_A\) :
\[
\rg(A)=\rg(u_A)=\dim(\Img(u_A)).
\]
\end{definition}

\begin{remarque}
Ainsi, le rang d’une matrice n’est pas un nouveau concept indépendant : c’est le rang de l’application linéaire qu’elle représente.
\end{remarque}

\subsection{Interprétation par les colonnes}

\begin{proposition}
Le rang d’une matrice est la dimension du sous-espace engendré par ses colonnes.
\end{proposition}

\begin{proof}
Écrivons les colonnes de \(A\) sous la forme
\[
A=
\begin{pmatrix}
| & & |\\
C_1 & \cdots & C_n\\
| & & |
\end{pmatrix}.
\]
Pour tout vecteur-colonne
\[
X=
\begin{pmatrix}
x_1\\
\vdots\\
x_n
\end{pmatrix},
\]
on a
\[
AX=x_1C_1+\cdots+x_nC_n.
\]
L’image de \(u_A\) est donc exactement l’espace engendré par \(C_1,\dots,C_n\). Ainsi
\[
\rg(A)=\dim(\Vect(C_1,\dots,C_n)).
\]
\end{proof}

\begin{corollaire}
Le rang d’une matrice est le rang de sa famille de colonnes.
\end{corollaire}

\subsection{Interprétation par les lignes}

\begin{theoreme}
Le rang d’une matrice est aussi la dimension du sous-espace engendré par ses lignes.
\end{theoreme}

\begin{remarque}
Autrement dit, le rang-colonne et le rang-ligne coïncident.
\end{remarque}

\begin{remarque}
Dans un premier cours, on peut l’admettre après avoir prouvé l’invariance du rang par opérations élémentaires et utilisé la forme échelonnée. Dans un poly plus avancé, on peut en donner une preuve complète via les matrices équivalentes.
\end{remarque}

\begin{corollaire}
Pour toute matrice \(A\),
\[
\rg(A)=\rg(A^T).
\]
\end{corollaire}

\section{Rang d’une famille et rang d’une matrice associée}

\begin{proposition}
Soit \((u_1,\dots,u_p)\) une famille de vecteurs de \(\K^n\).  
Si l’on forme la matrice
\[
A=
\begin{pmatrix}
| & & |\\
u_1 & \cdots & u_p\\
| & & |
\end{pmatrix},
\]
alors
\[
\rg(u_1,\dots,u_p)=\rg(A).
\]
\end{proposition}

\begin{proof}
Le rang de \(A\) est la dimension du sous-espace engendré par ses colonnes, c’est-à-dire par les vecteurs \(u_1,\dots,u_p\).
\end{proof}

\begin{remarque}
Ce résultat permet de calculer le rang d’une famille de vecteurs avec les outils du calcul matriciel.
\end{remarque}

\section{Rang et matrices échelonnées}

\subsection{Invariance par opérations élémentaires}

\begin{proposition}
Les opérations élémentaires sur les lignes ne changent pas le rang d’une matrice.
\end{proposition}

\begin{proof}
Une opération élémentaire sur les lignes revient à multiplier la matrice à gauche par une matrice élémentaire inversible \(E\). On obtient
\[
A'=EA.
\]
Comme \(E\) est inversible, l’application associée à \(A'\) a même rang que celle associée à \(A\). Donc
\[
\rg(A')=\rg(A).
\]
\end{proof}

\begin{remarque}
De même, les opérations élémentaires sur les colonnes ne changent pas le rang.
\end{remarque}

\subsection{Rang d’une matrice échelonnée}

\begin{proposition}
Le rang d’une matrice échelonnée est égal au nombre de pivots, c’est-à-dire au nombre de lignes non nulles.
\end{proposition}

\begin{proof}
Dans une matrice échelonnée, les lignes non nulles sont libres. Elles engendrent l’espace des lignes. Donc la dimension de l’espace des lignes est exactement le nombre de lignes non nulles. Comme le rang-ligne est égal au rang de la matrice, on obtient le résultat.
\end{proof}

\begin{corollaire}
Pour calculer le rang d’une matrice, on peut l’échelonner puis compter les lignes non nulles.
\end{corollaire}

\section{Calcul pratique du rang}

\begin{methode}
Pour calculer le rang d’une matrice :
\begin{enumerate}
    \item on effectue des opérations élémentaires sur les lignes ;
    \item on met la matrice sous forme échelonnée ;
    \item on compte les pivots ou les lignes non nulles.
\end{enumerate}
\end{methode}

\begin{exemple}
Calculons le rang de
\[
A=
\begin{pmatrix}
1 & 2 & 1\\
2 & 4 & 2\\
1 & 1 & 0
\end{pmatrix}.
\]
On effectue :
\[
L_2\leftarrow L_2-2L_1,
\qquad
L_3\leftarrow L_3-L_1.
\]
On obtient
\[
\begin{pmatrix}
1 & 2 & 1\\
0 & 0 & 0\\
0 & -1 & -1
\end{pmatrix}.
\]
Puis on échange \(L_2\) et \(L_3\) :
\[
\begin{pmatrix}
1 & 2 & 1\\
0 & -1 & -1\\
0 & 0 & 0
\end{pmatrix}.
\]
Il y a \(2\) lignes non nulles, donc
\[
\rg(A)=2.
\]
\end{exemple}

\section{Rang et systèmes linéaires}

\subsection{Système homogène}

\begin{proposition}
Pour \(A\in M_{p,n}(\K)\), l’ensemble des solutions du système homogène
\[
AX=0
\]
est le noyau de l’application linéaire \(u_A\). Ainsi
\[
\dim(\Ker(u_A))=n-\rg(A).
\]
\end{proposition}

\begin{proof}
Le système homogène
\[
AX=0
\]
décrit exactement les vecteurs \(X\in \K^n\) tels que
\[
u_A(X)=0.
\]
Donc son ensemble de solutions est \(\Ker(u_A)\). Le théorème du rang appliqué à \(u_A:\K^n\to\K^p\) donne
\[
n=\dim(\Ker(u_A))+\rg(A),
\]
d’où
\[
\dim(\Ker(u_A))=n-\rg(A).
\]
\end{proof}

\begin{remarque}
Cette formule donne le nombre de degrés de liberté d’un système homogène.
\end{remarque}

\subsection{Cas carré}

\begin{corollaire}
Si \(A\in M_n(\K)\), alors le système homogène
\[
AX=0
\]
admet uniquement la solution nulle si et seulement si
\[
\rg(A)=n.
\]
\end{corollaire}

\begin{proof}
Le système homogène n’a que la solution nulle si et seulement si
\[
\Ker(u_A)=\{0\},
\]
c’est-à-dire
\[
\dim(\Ker(u_A))=0.
\]
Puis on utilise
\[
\dim(\Ker(u_A))=n-\rg(A).
\]
\end{proof}

\section{Rang maximal et inversibilité}

\subsection{Matrices carrées}

\begin{theoreme}
Soit \(A\in M_n(\K)\). Les propriétés suivantes sont équivalentes :
\begin{itemize}
    \item \(\rg(A)=n\) ;
    \item les colonnes de \(A\) forment une base de \(\K^n\) ;
    \item les lignes de \(A\) forment une base de \(\K^n\) ;
    \item \(u_A\) est injective ;
    \item \(u_A\) est surjective ;
    \item \(u_A\) est bijective ;
    \item \(A\) est inversible.
\end{itemize}
\end{theoreme}

\begin{proof}
Comme \(u_A:\K^n\to\K^n\), le fait que \(\rg(A)=n\) équivaut à
\[
\dim(\Img(u_A))=n,
\]
donc à
\[
\Img(u_A)=\K^n.
\]
Ainsi \(u_A\) est surjective. En dimension finie égale, surjectivité et injectivité sont équivalentes.

La bijectivité de \(u_A\) équivaut à l’inversibilité de la matrice \(A\).

Enfin, les colonnes forment une base de \(\K^n\) si et seulement si elles engendrent \(\K^n\) avec \(n\) vecteurs, ce qui équivaut à un rang \(n\).
\end{proof}

\subsection{Cas général}

\begin{proposition}
Pour toute matrice
\[
A\in M_{p,n}(\K),
\]
on a
\[
\rg(A)\le \min(p,n).
\]
\end{proposition}

\begin{proof}
Le rang est la dimension d’un sous-espace de \(\K^p\), donc
\[
\rg(A)\le p.
\]
C’est aussi le rang d’une famille de \(n\) colonnes, donc
\[
\rg(A)\le n.
\]
\end{proof}

\section{Rang et composition}

\begin{proposition}
Soient
\[
u:E\to F,
\qquad
v:F\to G
\]
linéaires entre espaces de dimension finie. Alors
\[
\rg(v\circ u)\le \rg(u)
\qquad \text{et} \qquad
\rg(v\circ u)\le \rg(v).
\]
\end{proposition}

\begin{proof}
On a
\[
\Img(v\circ u)=v(\Img(u)).
\]
Donc \(\Img(v\circ u)\) est l’image par \(v\) d’un sous-espace de \(\Img(u)\), donc sa dimension ne peut pas dépasser celle de \(\Img(u)\). Ainsi
\[
\rg(v\circ u)\le \rg(u).
\]

Par ailleurs,
\[
\Img(v\circ u)\subset \Img(v),
\]
donc
\[
\rg(v\circ u)\le \rg(v).
\]
\end{proof}

\begin{corollaire}
Pour des matrices \(A\) et \(B\) de tailles compatibles,
\[
\rg(AB)\le \min(\rg(A),\rg(B)).
\]
\end{corollaire}

\section{Rang d’une somme}

\begin{proposition}
Soient \(u,v:E\to F\) deux applications linéaires. Alors
\[
\Img(u+v)\subset \Img(u)+\Img(v).
\]
\end{proposition}

\begin{proof}
Pour tout \(x\in E\),
\[
(u+v)(x)=u(x)+v(x)\in \Img(u)+\Img(v).
\]
\end{proof}

\begin{corollaire}
On a
\[
\rg(u+v)\le \rg(u)+\rg(v).
\]
\end{corollaire}

\begin{proof}
Comme
\[
\Img(u+v)\subset \Img(u)+\Img(v),
\]
on obtient
\[
\dim(\Img(u+v))\le \dim(\Img(u)+\Img(v)).
\]
Or
\[
\dim(\Img(u)+\Img(v))\le \dim(\Img(u))+\dim(\Img(v)).
\]
Donc
\[
\rg(u+v)\le \rg(u)+\rg(v).
\]
\end{proof}

\section{Rang des formes linéaires}

\begin{definition}
Une application linéaire
\[
f:E\to \K
\]
est appelée une \textbf{forme linéaire}.
\end{definition}

\begin{proposition}
Soit \(E\) un espace vectoriel de dimension finie \(n\), et \(f:E\to\K\) une forme linéaire non nulle. Alors
\[
\rg(f)=1
\qquad \text{et} \qquad
\dim(\Ker(f))=n-1.
\]
\end{proposition}

\begin{proof}
Comme \(f\neq 0\), son image n’est pas réduite à \(\{0\}\). Or \(\Img(f)\) est un sous-espace de \(\K\), donc
\[
\Img(f)=\K.
\]
Ainsi
\[
\rg(f)=1.
\]
Le théorème du rang donne alors
\[
n=\dim(\Ker(f))+1.
\]
\end{proof}

\begin{corollaire}
Le noyau d’une forme linéaire non nulle est un hyperplan.
\end{corollaire}

\section{Rang et lecture matricielle}

\subsection{Lecture colonne}

\begin{proposition}
Le rang d’une matrice est le nombre maximal de colonnes linéairement indépendantes.
\end{proposition}

\begin{proof}
Le rang est le rang de la famille de colonnes, donc la dimension du sous-espace qu’elles engendrent. C’est exactement le nombre maximal de colonnes libres.
\end{proof}

\subsection{Lecture ligne}

\begin{proposition}
Le rang d’une matrice est aussi le nombre maximal de lignes linéairement indépendantes.
\end{proposition}

\begin{proof}
Par égalité du rang-colonne et du rang-ligne.
\end{proof}

\section{Méthodes pratiques}

\begin{methode}
Pour calculer le rang d’une famille de vecteurs :
\begin{enumerate}
    \item on place les vecteurs en colonnes dans une matrice ;
    \item on échelonne ;
    \item on compte les pivots.
\end{enumerate}
\end{methode}

\begin{methode}
Pour calculer le rang d’une application linéaire :
\begin{enumerate}
    \item soit on calcule l’image d’une base ;
    \item soit on construit sa matrice et on en calcule le rang.
\end{enumerate}
\end{methode}

\begin{methode}
Pour exploiter le théorème du rang :
\begin{enumerate}
    \item calculer le noyau ou l’image ;
    \item puis utiliser
    \[
    \dim(E)=\dim(\Ker(u))+\rg(u).
    \]
\end{enumerate}
\end{methode}

\section{Erreurs classiques}

\begin{itemize}
    \item Confondre le rang avec le nombre de vecteurs d’une famille.
    \item Croire qu’une famille de \(p\) vecteurs a toujours le rang \(p\).
    \item Oublier que le rang d’une matrice se lit sur ses colonnes ou ses lignes, pas sur ses coefficients.
    \item Oublier l’invariance du rang par opérations élémentaires.
    \item Confondre
    \[
    \dim(\Ker(u))
    \quad \text{et} \quad
    \rg(u).
    \]
    \item Croire qu’une matrice carrée est toujours de rang maximal.
    \item Ne pas relier le rang d’un système homogène au nombre de paramètres libres.
\end{itemize}

\section{Résumé du chapitre}

\begin{itemize}
    \item Le rang d’une famille est la dimension du sous-espace qu’elle engendre.
    \item Le rang d’une application linéaire est la dimension de son image.
    \item Le rang d’une matrice est le rang de l’application linéaire associée.
    \item Le rang d’une matrice est la dimension de l’espace engendré par ses colonnes.
    \item Le théorème du rang donne
    \[
    \dim(E)=\dim(\Ker(u))+\rg(u).
    \]
    \item Pour \(A\in M_{p,n}(\K)\),
    \[
    \dim(\Ker(u_A))=n-\rg(A).
    \]
    \item Une matrice carrée d’ordre \(n\) est inversible si et seulement si
    \[
    \rg(A)=n.
    \]
    \item Le rang est invariant par opérations élémentaires.
    \item Le rang d’une matrice échelonnée est le nombre de pivots.
\end{itemize}

\section*{Exercices d’application immédiate}

\begin{enumerate}[label=\textbf{Exercice \arabic*.}]
    \item Déterminer le rang des familles suivantes :
    \[
    ((1,0),(0,1)),
    \qquad
    ((1,0),(2,0)),
    \qquad
    ((1,0,0),(0,1,0),(1,1,0)).
    \]

    \item Calculer le rang de
    \[
    A=
    \begin{pmatrix}
    1 & 2\\
    2 & 4
    \end{pmatrix}.
    \]

    \item Calculer le rang de
    \[
    A=
    \begin{pmatrix}
    1 & 2 & 1\\
    2 & 4 & 2\\
    1 & 1 & 0
    \end{pmatrix}.
    \]

    \item Soit
    \[
    u:\R^3\to\R^2,\qquad u(x,y,z)=(x+y,y+z).
    \]
    Déterminer \(\rg(u)\) puis \(\dim(\Ker(u))\).

    \item Montrer qu’une famille de \(p\) vecteurs est libre si et seulement si son rang vaut \(p\).

    \item Montrer qu’une matrice carrée est inversible si et seulement si son rang est maximal.

    \item Soit
    \[
    A\in M_{3,4}(\R)
    \]
    de rang \(2\). Déterminer la dimension de l’ensemble des solutions de
    \[
    AX=0.
    \]

    \item Soient
    \[
    u:E\to F,\qquad v:F\to G
    \]
    linéaires. Montrer que
    \[
    \rg(v\circ u)\le \min(\rg(u),\rg(v)).
    \]

    \item Soit \(f:E\to\K\) une forme linéaire non nulle, avec \(\dim(E)=n\). Montrer que \(\Ker(f)\) est un hyperplan.

    \item Soit
    \[
    A=
    \begin{pmatrix}
    1 & 0 & 1\\
    0 & 1 & 1
    \end{pmatrix}.
    \]
    Déterminer \(\rg(A)\) puis décrire \(\Ker(A)\).
\end{enumerate}

\section{Compléments sur le rang des matrices}

\subsection{Matrices extraites et mineurs}

\begin{definition}
On appelle \textbf{matrice extraite} d’une matrice \(A\) toute matrice obtenue en ne gardant que certaines lignes et certaines colonnes de \(A\).
Un \textbf{mineur d’ordre \(r\)} est le déterminant d’une matrice extraite carrée de taille \(r\).
\end{definition}

\begin{theoreme}
Pour toute matrice \(A\), on a \(\rg(A)=r\) si et seulement si il existe un mineur d’ordre \(r\) non nul et tous les mineurs d’ordre strictement supérieur sont nuls.
\end{theoreme}

\subsection{Rang de la transposée}

\begin{proposition}
Pour toute matrice \(A\),
\[
\rg(A^\top)=\rg(A).
\]
\end{proposition}

\begin{proof}
Les mineurs d’une matrice et de sa transposée ont les mêmes déterminants ; la caractérisation du rang par les mineurs donne alors l’égalité des rangs.
\end{proof}

\end{document}
